% Part: first-order-logic 
% Chapter: axiomatic-deduction 
% Section: proving-things-quant

\documentclass[../../../include/open-logic-section]{subfiles}

\begin{document}

\olfileid[hi]{fol}{axd}{prq}

\olsection{परिमाणकों सहित \usetoken{P}{derivation}}

\begin{ex}
$(\lforall[x][!A(x)] \land
\lforall[y][!B(y)]) \lif \lforall[x][(!A(x) \land !B(x))]$ की
!!a{derivation} दें।

पहले ध्यान दें कि
\begin{align*}
  (\lforall[x][!A(x)] \land \lforall[y][!B(y)]) & \lif \lforall[x][!A(x)]\\
  \intertext{\olref[prp]{ax:land1} का दृष्टान्त है, और}
  \lforall[x][!A(x)] & \lif !A(a) \\
  \intertext{\olref[qua]{ax:q1} का दृष्टान्त है। अतः \olref[pro]{prop:chain} से हम जानते हैं कि}
  (\lforall[x][!A(x)] \land \lforall[y][!B(y)]) & \lif !A(a) \\
  \intertext{!!{derivable} है। इसी प्रकार, क्योंकि}
  (\lforall[x][!A(x)] \land \lforall[y][!B(y)]) & \lif \lforall[y][!B(y)] \qquad\text{और}\\
    \lforall[y][!B(y)] & \lif !B(a)\\
    \intertext{क्रमशः \olref[prp]{ax:land2} और \olref[qua]{ax:q1} के दृष्टान्त हैं, इसलिए}
    (\lforall[x][!A(x)] \land \lforall[y][!B(y)]) & \lif !B(a)\\
    \intertext{\olref[pro]{prop:chain} से व्युत्पन्न है। \olref[prp]{ax:land3} के उपयुक्त दृष्टान्त और~\MP{} के दो प्रयोगों से मिलता है कि}
    (\lforall[x][!A(x)] \land \lforall[y][!B(y)]) & \lif (!A(a) \land !B(a))\\
    \intertext{व्युत्पन्न है। अब \QR{} लगाकर मिलता है}
(\lforall[x][!A(x)] \land \lforall[y][!B(y)]) & \lif \lforall[x][(!A(x) \land !B(x))].
\end{align*}
\end{ex}

\end{document}
